\begin{frame}{SVM: Khác gì giữa Linear và Kernel?}
    \begin{columns}
        \begin{column}{0.42\textwidth}
        \textbf{Linear SVM}:\\
            - Tìm một đường thẳng (hoặc siêu phẳng) để phân chia dữ liệu\\
            - Hoạt động trực tiếp trên dữ liệu gốc $x$\\
            - Phù hợp khi dữ liệu \textbf{có thể tách tuyến tính}
        \end{column}

        \begin{column}{0.57\textwidth}
        \textbf{Kernel SVM}:\\
            - Dùng \textbf{kernel} để biến dữ liệu sang không gian khác\\
            - Giúp tạo ra \textbf{đường phân chia cong} (phi tuyến)\\
            - Phù hợp khi dữ liệu \textbf{không tách được bằng đường thẳng}
        \end{column}
    \end{columns}

    \begin{alertblock}{Vì sao cần Kernel?}
        Thay vì biến đổi dữ liệu sang không gian nhiều chiều (rất tốn chi phí),
        ta chỉ cần tính \textbf{độ giống nhau} giữa hai điểm bằng hàm $K(x_i, x_j)$.\\
        Kernel Trick giúp:
        \begin{itemize}
            \item Xử lý được dữ liệu phi tuyến
            \item Không cần thực sự tính toán trong không gian cao chiều
            \item Giữ chi phí tính toán thấp hơn
        \end{itemize}
        % Đã xóa thẻ \end{itemize} bị thừa ở đây
    \end{alertblock}

    \vspace{0.2cm}

    \textbf{Hàm quyết định:}
    \[
        f(x) = \operatorname{sign}\left(\sum_i \alpha_i y_i K(x_i, x) + b\right)
    \]
\end{frame}

\begin{frame}{RBF Kernel: điểm khác biệt chính với Linear SVM}
                  Kernel được sử dụng phổ biến nhất là RBF kernel:
                \[
                    K(x_i, x_j) = \exp\left(-\frac{\|x_i - x_j\|^2}{2\sigma^2}\right)
                \]
                hoặc viết lại:
                \[
                    K(x_i, x_j) = \exp\left(-\gamma \|x_i - x_j\|^2\right)
                \]
    \begin{columns}[T]
        \begin{column}{0.48\textwidth}
            \begin{alertblock}{Ý nghĩa hình học}
                \begin{itemize}
                    \item Hàm Gaussian tâm ở mỗi điểm dữ liệu $x_i$
                    \item Độ tương tự giảm theo khoảng cách Euclid bình phương
                    \item $\sigma$ điều khiển độ rộng kernel
                \end{itemize}
            \end{alertblock}
        \end{column}

        \begin{column}{0.48\textwidth}

            \begin{block}{Lợi ích RBF}
                \begin{itemize}
                    \item Linh hoạt cho các bài toán phi tuyến
                    \item Dễ cài đặt và tối chỉnh
                    \item Hoạt động tốt với dữ liệu nhiều chiều
                \end{itemize}
                \vspace{0.5cm}
            \end{block}
        \end{column}
    \end{columns}

    \hfill{\scriptsize \cite{duchesnay2021statsml}}
\end{frame}

\begin{frame}[shrink=10]{SVM: Đánh giá}
            \begin{columns}[T]
                \begin{column}{0.48\textwidth}
                    \textbf{Ưu điểm:}
                    \begin{itemize}
                        \item Hiệu quả với dữ liệu nhiều chiều, biên phân tách rõ
                        \item Tổng quát hóa tốt nhờ tối ưu margin
                        \item Kết hợp kernel để xử lý quan hệ phi tuyến
                    \end{itemize}
                \end{column}
                \begin{column}{0.48\textwidth}
                    \textbf{Nhược điểm:}
                    \begin{itemize}
                        \item Nhạy với tham số $C$, $\gamma$ và lựa chọn kernel
                        \item Chi phí huấn luyện cao khi số mẫu lớn
                        \item Khó diễn giải trực quan hơn một số mô hình tuyến tính
                    \end{itemize}
                \end{column}
            \end{columns}

            \begin{alertblock}{Kết quả mẫu (Breast Cancer)}
                Sử dụng RBF kernel + StandardScaler:
                \begin{itemize}
                    \item \textbf{bAcc}: 0.97 (Balanced Accuracy)
                    \item \textbf{AUC}: 0.99 (công thức)
                    \item \textbf{AUC}: 0.99 (xác suất)
                \end{itemize}
            \end{alertblock}

    \hfill{\scriptsize \cite{duchesnay2021statsml}}
\end{frame}
